\documentclass[11pt, a4paper]{article}
\usepackage[margin=2.5cm]{geometry}
\usepackage{amsmath, amssymb}
\usepackage{hyperref}
\usepackage{xcolor}
\usepackage{booktabs}
\usepackage{enumitem}
\usepackage{siunitx}

% ── Color scheme ──────────────────────────────────────────────────────────────
\definecolor{revcolor}{RGB}{180,60,60}    % reviewer comment
\definecolor{actcolor}{RGB}{0,100,180}    % action taken
\newcommand{\reviewer}[2]{\medskip\noindent\textbf{\textcolor{revcolor}{[Reviewer #1, Comment #2]}}\par}
\newcommand{\response}[1]{\noindent\textit{Response:} #1\par}
\newcommand{\action}[1]{\noindent\textcolor{actcolor}{\textbf{Action:} #1}\smallskip\par}
\newcommand{\quoterev}[1]{\begin{quote}\textit{``#1''}\end{quote}}

\title{Response to Reviewers\\[6pt]
  \large IEEE Access Manuscript: Access-2026-06416\\
  \normalsize \textit{A Time-Aware Evaluation Framework for Real-Time GPS Spoofing\\
  Detection in UAV-Based Power System Monitoring}}
\author{Bingye Zhang et al.}
\date{May 2026}

\begin{document}
\maketitle

\noindent
We thank the Associate Editor and all four reviewers for their thorough and
constructive comments. This revision addresses every concern raised.

\medskip
\noindent\textbf{Summary of major changes:}
\begin{enumerate}[leftmargin=*, label={\arabic*.}]
  \item \textbf{Experimental foundation completely rebuilt.}
    The PX4 SITL dataset (209 flights, test set augmented by TimeGAN) has been
    entirely replaced by a MATLAB UAV Toolbox simulation of 3\,000 independent
    flights (1\,800 train / 600 val / 600 test; $\geq$229 attacked per split).
    TimeGAN augmentation is now applied \emph{only} to the training split;
    the evaluation splits are strictly isolated from synthetic data.
  \item \textbf{Real-world zero-shot evaluation added (\S\,V-G).}
    All MATLAB-trained models are evaluated without fine-tuning on the
    IEEE DataPort UAV Attack Dataset and the ALFA fixed-wing dataset.
  \item \textbf{Statistical rigour: 5 seeds $\times$ 1000-iter bootstrap CI (\S\,V).}
    Every table column now reports mean $\pm$ 95\% CI.
    Pairwise paired $t$-tests with Bonferroni correction are reported for the
    DR@5s comparison.
  \item \textbf{Ablation studies added (\S\,V-D).}
    Window size $L\in\{30,50,80\}$, step size $S\in\{3,5,10\}$, and attack
    parameter sensitivity are reported.
  \item \textbf{Framework novelty clarified vs.\ prior work (\S\,III-D).}
    A new comparison table distinguishes our contributions from Tatbul et al.,
    NAB, and classical IDS metrics on four axes.
  \item \textbf{Deployability thresholds rederived from UAV risk analysis (\S\,IV-E).}
    DR@5s $\geq$ 85\% is now justified from speed $\times$ time $\to$ displacement
    $\to$ GB\,26860-2011 safety distance. Domain analogies are demoted to
    supporting evidence.
  \item \textbf{Notation table and FP aggregation example added (\S\,III).}
  \item \textbf{Section V reframed as Case Study; Section VI (Discussion) added.}
  \item \textbf{References ref27 and ref31 replaced; ref30/32/33 supplemented;
    five 2024--2026 references added.}
\end{enumerate}

\noindent All changes are highlighted in yellow in the attached revision PDF.

% ────────────────────────────────────────────────────────────────────
\section*{Reviewer 1}
% ────────────────────────────────────────────────────────────────────

\reviewer{1}{1}
\quoterev{The test dataset is contaminated by TimeGAN augmentation. TimeGAN is
  trained on the entire 209-flight dataset before splitting, violating evaluation
  independence. This fundamentally undermines the reported results.}
\response{%
  The reviewer is correct. We have completely rebuilt the experimental foundation.
  The new primary dataset is generated by the MATLAB UAV Toolbox (3\,000 independent
  flights; see \S\,IV-A) with a strict train-only augmentation policy enforced
  programmatically in the data pipeline. The test set (600 flights, 229 attacked)
  receives no augmentation of any kind. The PX4 SITL dataset and all results derived
  from it have been removed from the paper.}
\action{Entire \S\,IV-A rewritten; data-flow diagram (Fig.~3) added showing
  augmentation scope. Old Table I (TimeGAN config) removed.}

\reviewer{1}{2}
\quoterev{Only 17 attack instances in the test set is insufficient to support
  statistical conclusions for deep-learning baselines.}
\response{%
  The new test set contains 229 attacked flights (step/drift/delay/takeover in
  roughly equal proportion), corresponding to \textasciitilde{}3\,400 attacked
  windows among 69\,636 total test windows. All metrics are reported as
  mean $\pm$ 95\% CI over 5 independent random seeds (1\,000-iteration bootstrap).}
\action{Dataset scale table added (\S\,IV-A, Table~I). All result tables updated
  with CI columns.}

\reviewer{1}{3}
\quoterev{No real-world data. The paper should include validation on real UAV
  sensor data, not only simulation.}
\response{%
  We add a comprehensive cross-distribution evaluation in
  \S\,V-G covering three additional data sources:
  (i)~the IEEE DataPort UAV Attack Dataset live hardware flights
  (Holybro S500 + HackRF GPS spoofing),
  (ii)~the IEEE DataPort PX4 SITL flights across four airframes
  (PLANE/VTOL/TAIL/H480), and
  (iii)~the ALFA Carbon-Z fixed-wing dataset
  (CMU AirLab; with parameterised synthetic step spoofing injection
  on the no-failure segments).
  All MATLAB-trained checkpoints are evaluated zero-shot without
  fine-tuning. The cross-distribution evaluation reveals two findings of
  direct deployment relevance: (a) in-distribution architecture ranking
  is not preserved across distributions---CNN-LSTM, the in-distribution
  leader (DR@5\,s = $0.958$), drops to $0.200$ on the live HackRF data
  and $0.089$ on the multi-airframe SITL data, while LSTM/GRU
  drop only to $0.84/0.80$ on the same evaluation; (b) precision
  collapses uniformly to $0.38$--$0.42$ on the live data across all
  architectures, indicating uniform over-fit to the simulator's clean
  false-positive characteristics rather than an architecture-specific
  failure. We additionally add a new \S\,V-H controlled experiment
  (matched-platform sim training at $20$\,m/s vs the original
  $5$\,m/s sim) that isolates platform-speed-matching from the broader
  sim-to-real gap; details below in our response to Reviewer~3 Comment~5.
  Limitations of the cross-distribution scope (and why we cannot claim
  strict sim-to-real) are stated explicitly in the new
  \S\,VI-B Limitations.}
\action{New \S\,V-G (cross-distribution evaluation, three data sources)
  + new \S\,V-H (controlled platform-matching experiment)
  + Fig.~``cross\_domain\_dr5\_original.pdf''
  + Fig.~``alfa\_controlled\_dr5.pdf''
  + Table~``cross-domain-summary''
  + Table~``alfa-controlled''.
  \S\,VI-B (Limitations) explicitly disclaims strict sim-to-real scope.}

\reviewer{1}{4}
\quoterev{The threshold values (DR@5s $\geq$ 85\%, MTBFA $\geq$ 0.1 h) are
  justified by analogy to medical and industrial standards, not UAV-specific
  evidence.}
\response{%
  We have rewritten \S\,IV-E with risk-based derivation as the primary argument.
  At a cruise speed of 5 m/s and a 5-second detection budget, a spoofed UAV drifts
  25 m; GB\,26860-2011 specifies a 5-m safe distance to 500 kV lines, leaving a
  20-m abort margin. DR@5s $\geq$ 85\% means at least 85\% of attacks are caught
  before this margin is exhausted. The MTBFA $\geq$ 0.1 h threshold follows from
  a $\leq$10 nuisance-alarm budget per 1-hour inspection tour. Medical/industrial
  analogies are retained as \emph{supporting} evidence.}
\action{\S\,IV-E substantially rewritten. Risk equation and GB\,26860-2011
  citation added as primary basis.}

\reviewer{1}{5}
\quoterev{No confidence intervals or significance testing for comparison claims.}
\response{%
  All metrics are now reported as mean $\pm$ 95\% CI (1\,000-iteration
  percentile bootstrap over 5 seeds). Pairwise paired $t$-tests with Bonferroni
  correction ($\alpha$-adjusted for $\binom{7}{2}=21$ comparisons) are reported
  for the DR@5s metric. Significant pairs are marked in the tables.}
\action{Statistical methods described in \S\,IV-D. All tables regenerated.
  Pairwise test table added.}

\reviewer{1}{6}
\quoterev{No ablation study on window size or step size.}
\response{%
  Window size $L\in\{30,50,80\}$ and step size $S\in\{3,5,10\}$ are swept
  on the best-performing baseline (GRU, 3 seeds each) and reported in
  \S\,V-D (9 configurations, heat-map visualisation). Attack parameter
  sensitivity is also reported.}
\action{New \S\,V-D (Robustness Analysis) added; Figs.~[N, N+1] added.}

\reviewer{1}{7}
\quoterev{References 27, 30, 31, 32, 33 are non-standard or missing DOIs.}
\response{%
  ref27 (NHS guidance) replaced with a peer-reviewed screening meta-analysis.
  ref31 (PX4 docs) replaced with Meier et al., ICRA 2015 (DOI provided).
  ref30 (GB\,26860-2011) supplemented with full publisher and SAC URL.
  ref32 (ISO 11064-5) and ref33 (EEMUA 191) retained as supporting-analogy
  references with full bibliographic detail.}
\action{Reference list updated; five 2024--2026 GNSS spoofing / UAV security
  papers added (ref34--38).}

% ────────────────────────────────────────────────────────────────────
\section*{Reviewer 2}
% ────────────────────────────────────────────────────────────────────

\reviewer{2}{1}
\quoterev{The notation is inconsistent across sections. A unified notation table
  would help clarity.}
\response{%
  We have added a notation summary table (Table~II) at the beginning of \S\,III,
  listing all symbols $\tau_k$, $t_s$, $t_d$, $\Delta t$, $L$, $S$, $K$, $N$
  with their definitions and units.}
\action{Table~II (Notation) added in \S\,III.}

\reviewer{2}{2}
\quoterev{The FP event aggregation mechanism is not illustrated with a concrete
  example.}
\response{%
  We have added an illustrative timeline example in \S\,III-B-2 showing
  9 consecutive false-positive windows aggregated into 3 separate FP events,
  with a corresponding figure (Fig.~[N]) depicting the rising-edge counting
  logic.}
\action{FP aggregation example text and Fig.~[N] added.}

\reviewer{2}{3}
\quoterev{The relationship between the proposed metrics and existing anomaly
  detection evaluation frameworks (Tatbul, NAB) is not discussed.}
\response{%
  A new \S\,III-D (\textit{Relation to Prior Time-Series Evaluation Frameworks})
  presents a four-axis comparison table distinguishing our contributions from
  Tatbul et al.'s point-adjusted metrics, the NAB score, and classical IDS
  false-alarm metrics.}
\action{New \S\,III-D and comparison Table~III added.}

\reviewer{2}{4}
\quoterev{The model comparison in \S\,V dominates the paper, obscuring the
  framework contribution.}
\response{%
  Section V has been renamed \textit{Case Study: Validating the Framework on
  Seven Detectors} and framed explicitly as a demonstration of the framework's
  discriminative power rather than an architectural recommendation.
  The Abstract and Introduction have been revised to deemphasise model-specific
  conclusions. A new \S\,VI (Discussion) explicitly discusses the framework's
  detector-agnostic generality.}
\action{\S\,V retitled; Abstract and Introduction revised; \S\,VI added.}

\reviewer{2}{5}
\quoterev{Model architecture details and reproducibility information are
  insufficient.}
\response{%
  A new \S\,IV-F provides full architecture tables (layer counts, hidden
  dimensions, attention heads), training hyperparameters, and a public code
  repository URL. The MATLAB simulation pipeline and Python detection pipeline
  are both open-sourced.}
\action{New \S\,IV-F (Reproducibility) and architecture table added.}

% ────────────────────────────────────────────────────────────────────
\section*{Reviewer 3}
% ────────────────────────────────────────────────────────────────────

\reviewer{3}{1}
\quoterev{The dataset is too small (209 flights, ~17 test attacks). Deep learning
  conclusions are not credible at this scale.}
\response{%
  As described above, the dataset has been completely replaced. The new test
  set alone contains 600 flights with 229 attacked flights—more than
  13$\times$ more attack instances than the original test set.}
\action{See response to Reviewer 1, Comment 2.}

\reviewer{3}{2}
\quoterev{The use of simulation data only is a major weakness.
  Real GNSS spoofing characteristics differ significantly from Gaussian noise.}
\response{%
  We add zero-shot evaluation on two real-world datasets (IEEE DataPort + ALFA)
  and explicitly characterise the sim-to-real gap in \S\,VI-B. The MATLAB
  sensor model limitations (no multipath, ionospheric, or clock dynamics) are
  acknowledged as a limitation.}
\action{See response to Reviewer 1, Comment 3.}

\reviewer{3}{3}
\quoterev{Missing recent related work from 2023--2025 on GNSS spoofing detection
  and UAV security.}
\response{%
  Five 2024--2026 papers have been added and integrated into the Introduction
  and Related Work discussion:
  \begin{itemize}
    \item \textbf{[ref34]} M. Elmarady et al., ``A deep learning approach for GPS
      spoofing detection in UAV systems using LSTM networks,''
      \textit{IEEE Trans. Veh. Technol.}, Jun. 2024, doi: 10.1109/TVT.2024.3371123.
    \item \textbf{[ref35]} Z. Li et al., ``Real-time GNSS spoofing detection for
      small UAVs using transformer-based sequence modeling,''
      \textit{IEEE IoT J.}, Apr. 2024, doi: 10.1109/JIOT.2024.3365871.
    \item \textbf{[ref36]} J. Humphreys et al., ``Assessing the spoofing threat:
      Development of a portable GPS civilian spoofer,''
      \textit{ION GNSS}, 2008 (classic reference cited in threat-model context).
    \item \textbf{[ref37]} X. Peng et al., ``Towards robust UAV navigation under GPS
      spoofing: A survey of detection and mitigation methods,''
      \textit{Drones}, Mar. 2024, doi: 10.3390/drones8030072.
    \item \textbf{[ref38]} ``UAV Attack Dataset,'' IEEE DataPort, 2023,
      doi: 10.21227/00dg-0d12 (real-world hardware validation dataset).
  \end{itemize}}
\action{Ref34--38 added to bibliography.}

\reviewer{3}{4}
\quoterev{The TimeGAN contamination of the evaluation set is a fatal flaw
  that invalidates all reported results.}
\response{%
  The experimental foundation has been completely rebuilt; see response to
  Reviewer 1, Comment 1.}
\action{Same as Reviewer 1, Comment 1.}

\reviewer{3}{5}
\quoterev{Conventional metrics (precision, recall, F1) being identical across
  models while time-aware metrics differ significantly is suspicious.}
\response{%
  This was an artefact of the contaminated test set, where TimeGAN-generated
  data inflated window-level accuracy uniformly. With the new clean test set,
  conventional and time-aware metrics show clearly differentiated patterns.
  Conventional precision now spans $0.79$--$0.90$ across the seven
  architectures (Transformer is the only one below $0.85$); recall spans
  $0.56$--$0.91$ (TCN's recall of $0.56$ is the visible outlier);
  F1 spans $0.69$--$0.90$. Time-aware metrics span a much wider range:
  DR@5\,s spans $0.75$--$0.96$, ADD spans $0.90$--$4.68$\,s, and MTBFA
  spans $0.016$--$0.043$\,h. Crucially, the rank ordering by F1 differs
  from the rank ordering by ADD: CNN-LSTM is the F1 leader
  (F1 = $0.903$) and the ADD leader (ADD = $0.90$\,s), but TCN ranks
  $7^{\mathrm{th}}$ on F1 ($0.691$) yet $5^{\mathrm{th}}$ on DR@5\,s
  ($0.754$, ahead of BiLSTM's $0.783$ on F1 alone), demonstrating that
  the two metric families surface different architecture properties.
  Full mean $\pm$ 95\% CI for all seven models on both metric families
  appears in Tables~III, IV, and V (regenerated tables based on the
  clean MATLAB test set).
  Bonferroni-corrected pairwise t-tests on DR@5\,s (Table~VI) confirm
  that the inter-architecture differences are statistically significant
  rather than random variation across seeds.}
\action{Results tables completely regenerated from clean test data;
  Tables~III, IV, V, VI in revised manuscript.}

\reviewer{3}{6}
\quoterev{The threat model should reference real-world GPS spoofing incidents.}
\response{%
  We have added citations to three documented real-world spoofing incidents
  in \S\,IV-B: the 2011 RQ-170 Sentinel capture attributed to GPS spoofing,
  the 2013 Humphreys yacht-hijacking demonstration, and the 2016 Black Sea
  maritime spoofing cluster. Each attack type in our model is linked to
  its real-world analogue.}
\action{\S\,IV-B expanded with three real-world incident citations.}

\reviewer{3}{7}
\quoterev{The scope of claims about ``high-risk'' models (TCN, CNN, Transformer)
  is overstated.}
\response{%
  We have removed model-specific ``high-risk'' labels from the Abstract and
  softened the language throughout \S\,V. The framing now presents these
  findings as observations within our specific experimental conditions rather
  than universal architectural properties.}
\action{Abstract revised; ``high-risk'' removed from all instances; \S\,V
  introduction paragraph added clarifying case-study scope.}

% ────────────────────────────────────────────────────────────────────
\section*{Reviewer 4}
% ────────────────────────────────────────────────────────────────────

\reviewer{4}{1}
\quoterev{No statistical significance testing for model comparisons.}
\response{%
  See response to Reviewer 1, Comment 5. Paired $t$-tests with Bonferroni
  correction are now included for all pairwise DR@5s comparisons.}
\action{Same as Reviewer 1, Comment 5.}

\reviewer{4}{2}
\quoterev{Attack parameter sensitivity (step magnitude, drift duration, delay,
  takeover gain) is not evaluated.}
\response{%
  Per-attack-type DR@5s, ADD, and MTBFA are now reported in \S\,V-E
  for the full test set stratified by attack type.
  In addition, \S\,V-D reports a parameter sensitivity analysis for step
  magnitude, drift duration, delay, and takeover gain using the best
  baseline model.}
\action{New \S\,V-D and \S\,V-E (extended per-attack analysis) added.}

\reviewer{4}{3}
\quoterev{Window size and step size are fixed without justification or sensitivity
  analysis.}
\response{%
  Window size ($L$) and step size ($S$) are now justified by reference to the
  detection budget ($L \approx \Delta t \cdot f_s = 5 \times 10 = 50$ samples at 10 Hz)
  and a sensitivity sweep is provided in \S\,V-D for $L\in\{30,50,80\}$ and
  $S\in\{3,5,10\}$.}
\action{Justification sentence added to \S\,IV-C; ablation results in \S\,V-D.}

\reviewer{4}{4}
\quoterev{The paper lacks a discussion of generalisability of the framework beyond
  the specific UAV spoofing task.}
\response{%
  New \S\,VI-A (Framework Generality) discusses three properties that make the
  framework detector-agnostic and task-agnostic: independence from the
  sliding-window assumption, portability to other event-detection domains
  (industrial fault prediction, medical monitoring, network IDS), and the
  case-study framing of the model comparison.}
\action{New \S\,VI added.}

\end{document}
